%
% 第9章：参数量与计算量分析

\clearpage
\cleardoublepage

\chapter{参数量与计算量分析}

\section{模型效率的重要性}

理解模型效率对于以下方面至关重要：
\begin{itemize}
    \item 在资源受限设备上部署
    \item 训练成本估算
    \item 模型选择和架构设计
    \item 基准测试和比较
\end{itemize}

关键指标：
\begin{itemize}
    \item \textbf{参数量：} 模型大小和内存占用
    \item \textbf{FLOPs：} 计算复杂度
    \item \textbf{延迟：} 推理时间
    \item \textbf{吞吐量：} 每秒样本数
\end{itemize}

\section{参数计数}

参数是网络中的可学习权重和偏置。

\subsection{全连接层}

对于输入维度 $d_{in}$ 和输出维度 $d_{out}$：
\begin{align}
\text{权重} &= d_{in} \times d_{out} \\
\text{偏置} &= d_{out} \\
\text{总计} &= d_{in} \times d_{out} + d_{out}
\end{align}

示例：784到256的全连接层：
\begin{equation}
784 \times 256 + 256 = 200,960 \text{ 个参数}
\end{equation}

\subsection{卷积层}

对于输入通道 $C_{in}$，输出通道 $C_{out}$，卷积核大小 $K$：
\begin{align}
\text{权重} &= K \times K \times C_{in} \times C_{out} \\
\text{偏置} &= C_{out} \\
\text{总计} &= K^2 \times C_{in} \times C_{out} + C_{out}
\end{align}

\section{FLOPs（浮点运算次数）}

FLOPs衡量前向传播的计算成本。

\subsection{卷积层FLOPs}

对于输出特征图 $H_{out} \times W_{out}$：
\begin{align}
\text{每个输出像素FLOPs} &= K^2 \times C_{in} \times C_{out} \times 2 \\
\text{总FLOPs} &= H_{out} \times W_{out} \times K^2 \times C_{in} \times C_{out} \times 2
\end{align}

示例：$224\times224$ 输入，$3\times3$ 卷积，64入，128出：
\begin{align}
H_{out} \times W_{out} &= 224 \times 224 = 50,176 \\
\text{FLOPs} &= 50,176 \times 3 \times 3 \times 64 \times 128 \times 2 \approx 1.17 \text{ GFLOPs}
\end{align}

\section{实际示例}

\begin{table}[h]
    \centering
    \caption{AlexNet参数统计}
    \begin{tabular}{L{2cm} L{2.5cm} L{2cm} L{2cm}}
        \toprule
        \textbf{层} & \textbf{输出} & \textbf{参数} & \textbf{FLOPs} \\
        \midrule
        Conv1 & $96\times55\times55$ & 35K & 105M \\
        Conv2 & $256\times27\times27$ & 614K & 448M \\
        Conv3 & $384\times13\times13$ & 885K & 149M \\
        Conv4 & $384\times13\times13$ & 1.3M & 224M \\
        Conv5 & $256\times13\times13$ & 884K & 149M \\
        FC6 & 4096 & 37.7M & 37.7M \\
        FC7 & 4096 & 16.8M & 16.8M \\
        FC8 & 1000 & 4.1M & 4.1M \\
        \midrule
        总计 & - & 约60M & 约1.1B \\
        \bottomrule
    \end{tabular}
\end{table}

\section{迁移学习与效率}

迁移学习可以显著减少所需数据和训练成本。

\begin{properties}[迁移学习优势]
\begin{itemize}
    \item 更低的训练成本
    \item 可使用更小的数据集
    \item 收敛更快
    \item 通常性能更好
\end{itemize}
\end{properties}

常见策略：
\begin{enumerate}
    \item \textbf{冻结：} 保持预训练权重固定
    \item \textbf{微调：} 以小学习率调整部分层
    \item \textbf{特征提取：} 使用预训练作为特征提取器
\end{enumerate}

\section{本章小结}

\begin{enumerate}
    \item 参数量决定模型大小和内存需求
    \item FLOPs衡量计算复杂度
    \item MobileNet等高效架构减少两者
    \item 迁移学习高效利用预训练模型
    \item 理解效率指导模型选择
\end{enumerate}
